
In this section we will study how the design parameters translate to
the actuation matrix $\mathbf{A}$ defined in the previous section. In
particular, we will consider how actuation limits translates into
maximum values of force and torque. We will then optimize these
parameters in order to maximize the upper bound of forces and torques,
over all directions, imposed by a bounded actuation.

\subsection{Problem statement}
\label{sec:problem-statement}

% Throughout this paper we will assume that the propellers are equally
% distributed radially, that is,
% \begin{equation}
%   \label{eq:thetas}
%   \theta_i = (i-1) \frac{\pi}{3}
% \end{equation}
% For the sake of symmetry we will further consider the rotation
% directions to be alternating, that is,
% \begin{equation}
%   \label{eq:ws}
%   w_i =
%   \begin{cases}
%     1 & i\text{ is odd,} \\
%     -1 & i\text{ is even.}
%   \end{cases}
% \end{equation}

We will start by considering that each actuation signal is
bounded between $-1$ and $1$, that is,
\begin{equation}
  \label{eq:ahcube}
  -1 \leq u_i \leq 1 \quad \text{for} \quad i=1,\ldots,6
\end{equation}
According to \eqref{eq:au}, this hypercube will map onto a
6-dimensional convex polyhedron\footnote{A convex polyhedron is an
  intersection of a finite number of half-spaces~\cite{jeter86}.} in
the $(\bar{F}, \bar{M})$ space.  Any other choice of bounds is
possible by appropriately scaling constants $K_1$ and $K_2$.  However,
it assumes that the maximum propeller trust is symmetric with respect
to the direction of rotation~\footnote{Propellers may have different
thrust capabilities along each direction of the rotation axis. The
approach presented here can be easily extended to account for that asymmetry.}.
The remaining parameters are the angles
$\{\phi_i\}$. We will base our analysis on the optimization of these
angles with respect to various criteria.

Our goal will be to find the configurations of angles $\{\phi_i\}$
that maximize the range of forces (and torques) over all
directions.
% Conventional multirotors have all propellers directed at
% the same direction, \textit{i.e.,} with $\{\phi_i=0\}$,, which
% maximizes the thrust along that direction. However, they are unable to
% produce any thrust component over an orthogonal direction. By having
% nonzero values of $\{\phi_i\}$
% we expect to gain thrust over any direction, at the
% expense of reducing the maximum thrust over some particular one.
Geometrically, this corresponds to changing $\{\phi_i\}$
such that a ball of nonzero radius can fit inside the 3-dimensional
convex polyhedron in the $\bar{F}$ space mapped by the actuation
hypercube in~\eqref{eq:ahcube}, while keeping zero torque,
$\bar{M}=0$. A similar reasoning applies to the torque space
$\bar{M}$, while keeping $\bar{F}=0$. 
% In sum, we will look for the values of $\{\phi_i\}$ that
% maximize these radii.

First, we will address the problem of computing the maximum force
along a given direction specified as a unit vector $\hat{e}$, while
maintaining a zero torque.  From~\eqref{eq:au}, and assuming that
$\mathbf{A}$ is full rank, we get
\begin{equation}
  \bar{u} = \mathbf{A}^{-1}
  \left(
    \begin{array}{c}
      \bar{F} \\
      \bar{M}
    \end{array}
  \right) =\mathbf{A}^{-1}
  \left(
    \begin{array}{c}
      F \hat{e} \\
      0
    \end{array}
  \right)
\end{equation}
where $F>0$ is the force magnitude. For what follows, it will be
convenient to express the $\mathbf{A}^{-1}$ matrix as blocks of three
dimensional row vectors:
\begin{equation}
  \label{eq:bici}
  \mathbf{A}^{-1} = \left[
    \begin{array}{cc}
      b_1^T & c_1^T \\
      \vdots & \vdots \\
      b_6^T & c_6^T \\
    \end{array}
    \right]
\end{equation}
where $b_i,c_i\in\mathbb{R}^3$. Then, the actuation of the $i$-th
propeller is given by
\begin{equation}
  u_i = F b_i^T \hat{e}
\end{equation}
Since $|u_i|\leq1$, we have $F|b_i^T\hat{e}|\leq1$, and thus $F$ has
this upper bound:
\begin{equation}
  F \leq \frac{1}{|b_i^T\hat{e}|}
\end{equation}
Since this inequality has to be satisfied for all propellers
$i=1,\ldots,6$, the maximum force $F^\mathrm{max}_{\hat{e}}$ is given
by the lowest of these upper bounds
\begin{equation}
  F^\mathrm{max}_{\hat{e}} = \min_i \frac{1}{|b_i^T\hat{e}|}
\end{equation}
This force is the maximum force along a given direction $\hat{e}$. The
maximum force attainable in any direction can be obtained by
minimising this force over all possible directions. Since
$|b_i^T\hat{e}|\leq\|b_i\|$, this minimum is given by
\begin{equation}
  \label{eq:fmax}
  F^\mathrm{max} = \min_i \frac{1}{\|b_i\|}
\end{equation}
The same reasoning can be applied to the torques: consider a torque
$\bar{M}=M\hat{e}$ along an arbitrary direction defined by $\hat{e}$,
the corresponding actuation with $\bar{F}=0$ is $u_i=M c_i^T \hat{e}$,
resulting in the following maximum torque along any direction:
\begin{equation}
  \label{eq:mmax}
  M^\mathrm{max} = \min_i \frac{1}{\|c_i\|}
\end{equation}

Now, these maximum forces and torque values depend on the design
parameters. In the following we will use an optimization approach to
find the values of these parameters that maximize the maximum force
and/or torque. We will consider the propellers to be equally
distributed radially, that is,
\begin{equation}
  \label{eq:thetas}
  \theta_i = (i-1) \frac{\pi}{3}
\end{equation}
and a fixed distance $d$, as well as the constants $K_1$ and $K_2$. All the
remaining parameters will be the unknown variables:
\begin{equation}
  \bar{\psi} = ( \phi_1, \ldots, \phi_6, w_1,\ldots, w_6 )^T
\end{equation}
with the feasibility domain defined by
\begin{equation}
  \Psi = \left\{ \bar{\psi} \::\: |\phi_{1,\ldots,6}| \leq
  \phi_{max}, w_{1,\ldots,6}\in\{-1,1\} \right\}
\end{equation}
where $\phi_{max}$ is the maximum allowed deviation from the
vertical. As a mechanical constraint to allow obstructionless air flow
we considered $\phi_{max}=\pi/3$ in this work.

We can restate the maximization of~\eqref{eq:fmax}
and~\eqref{eq:mmax} using the epigraph form~\cite{boyd04}, thus getting rid of the
maximization of a minimum, in exchange of the introduction of two new
variables, $p$ and $q$:
\begin{equation}
  \label{eq:minp}
  \begin{split}
    &\text{minimize}\: p \\
    &\text{subject to:} \\
    &\quad p \geq \|b_i\|^2, i=1,\ldots,6
  \end{split}
\end{equation}
for the force and
\begin{equation}
  \label{eq:minq}
  \begin{split}
    &\text{minimize}\: q \\
    &\text{subject to:} \\
    &\quad q \geq \|c_i\|^2, i=1,\ldots,6
  \end{split}
\end{equation}
for the torque, where $\{b_i\}$ and $\{c_i\}$ depend non-linearly on
the parameters $\bar{\psi}$ through the inverse of the actuation
matrix $\mathbf{A}$ as~\eqref{eq:bici}. In this form, the optimization
variables are augmented with the cost, that is,
$\bar{\psi}_p=(p, \bar{\psi})$ for the problem~\eqref{eq:minp} and
$\bar{\psi}_q=(q, \bar{\psi})$ for~\eqref{eq:minq}.  It can be readily
seen that these forms maximize~\eqref{eq:fmax} and~\eqref{eq:mmax},
where the resulting maximum forces and torques can be recovered using
$F^{max}=1/\sqrt{p}$ and $M^{max}=1/\sqrt{q}$.

\subsection{Multi-criteria optimization}
\label{sec:multi-crit-optim}

The design optimization aims at maximizing simultaneously two
objective functions, force and torque ranges, given bounded
actuation~\eqref{eq:ahcube}. To do so we will take a multi-criteria
optimization approach~\cite{statnikov95}, combining the
optimization problems~\eqref{eq:minp}
and~\eqref{eq:minq} into a single multi-criteria problem:
\begin{equation}
  \label{eq:multi}
  \begin{split}
    &\text{minimize}\:(p, q) \\
    &\text{subject to:} \\
    &\quad p \geq \|b_i\|^2, i=1,\ldots,6 \\
    &\quad q \geq \|c_i\|^2, i=1,\ldots,6
  \end{split}
\end{equation}
In this problem, the optimization variables are augmented with both
$p$ and $q$, $\bar{\psi}_{pq}=(p, q, \bar{\psi})$, and we have two
cost functions, say $J_1(\bar{\psi}_{pq})=p$ and
$J_2(\bar{\psi}_{pq})=q$. The solution of this multi-criteria
optimization problem is the set $\mathcal{P}$ of non-dominated
solutions, defined by: $\bar{\psi}_{pq}^0\in P$ if and only if there is no
$\bar{\psi}_{pq}\in\Psi$ such that
$J_i(\bar{\psi})\leq J_i(\bar{\psi}^0)$ for all $i\in\{1,2\}$ and
$J_i(\bar{\psi})< J_i(\bar{\psi}^0)$ for at least one
$i\in\{1,2\}$. This set is also called \emph{Pareto optimal
  set}~\cite{statnikov95}, a subset of the \emph{objective space}
defined by all~$(J_1(\bar{\psi}),J_2(\bar{\psi}))$ for
$\bar{\psi}\in\Psi$.

Apart from very simple cases, the Pareto optimal set is not trivial to
obtain exactly. Thus, we will make a pointwise approximation using the
Normally Boundary Intersection (NBI) method~\cite{das98}. This method
is guaranteed to obtain Pareto optimal points if the objective space
is convex. But it is still capable of obtaining points in ``sufficiently
concave'' parts of the objective space~\cite{das98}.

The first step of NBI is to obtain the minimizers of the each cost
function taken individually. These are also called \emph{shadow
  minima}. Let us start by considering the first minimization
problem~\eqref{eq:minp}, where $\bar{\psi}^*_p$ is the minimizer with
minimum cost $p^*$.  Then, this minimizer both minimizes $p$
in~\eqref{eq:multi} and, together with $q^0=\max_i\,\|c_i\|^2$, is a
non-dominated solution of~\eqref{eq:multi}, and thus belongs to its
Pareto optimal set. This results from the fact that this $q^0$ is the
smallest one that still satisfies the constrains of~\eqref{eq:multi}:
any $(p^*,q)$ with $q>q^0$ is dominated by $(p^*,q^0)$.  The same
reasoning can be applied to~\eqref{eq:minq}, resulting in the
minimizer $\bar{\psi}^*_q$, with minimum cost $q^*$, that together
with $p^0=\max_i\,\|b_i\|^2$ is also a non-dominated solution
of~\eqref{eq:multi}, thus also belonging to its Pareto optimal set.

On the $(p,q)$ space, these two non-dominated solutions corresponds
to two extremal points, $(p^*,q^0)$ and $(p^0,q^*)$, of the Pareto
optimal set: no feasible solutions exists neither to the left of
$p^*$ nor lower than $q^*$. \emph{The application of NBI to a two cost
function problem amounts to scanning along a straight line joining
$(p^*,q^0)$ and $(p^0,q^*)$, and then, for each point on this
straight line, to determine the single non-dominated solution along
the orthogonal direction.}

\begin{figure}
  \centering
  \includegraphics[width=0.6\linewidth]{nbi.pdf}
  \caption{Illustration of the NBI method: given a point defined by
    $\lambda\in[0;1]$, along the search line between $(p^*,q^0)$ and
    $(p^0,q^*)$, the optimization is done along the normal direction
    spawned by vector $\bar{n}$, resulting on the $NBI(\lambda)$
    intersection point.}
  \label{fig:nbi}
\end{figure}

Figure~\ref{fig:nbi} illustrates the NBI method. This straight line
can be parametrized by a $\lambda$ value ranging between 0 and 1,
resulting in $(1-\lambda)(p^*,q^0)+\lambda\,(p^0,q^*)$. An
orthogonal direction to this straight line is spawned by the vector
$\bar{n}=(q^0-q^*,p^0-p^*)$. Using again the epigraph form, but now
along this vector, we obtain the following constrained optimization
problem:
\begin{equation}
  \label{eq:nbi}
  \begin{aligned}
    &\text{minimize}\: t \\
    &\text{subject to:} \\
    &\quad (q^0-q^*)\,t + (1-\lambda) p^* + \lambda p^0 &\geq
    \|b_i\|^2 \\
    &\quad (p^0-p^*)\,t + \lambda q^* + (1-\lambda) q^0 &\geq
    \|c_i\|^2 \\
    &\qquad \text{for } i=1,\ldots,6
  \end{aligned}
\end{equation}
with the augmented vector $\bar{\psi}_t=(t, \bar{\psi})$ as
optimization variable. For a given $\lambda\in[0;1]$, the solution of
this optimization problem yields a minimizer $\bar{\psi}^*_t(\lambda)$
from which the corresponding point in the $(p,q)$ space is
\begin{equation}
  NBI(\lambda) = (p_\lambda, p_\lambda), \quad p_\lambda = \min_i \|b_i\|^2, \quad q_\lambda = \min_i \|c_i\|^2
\end{equation}
from which the maximum values of force and torque can be recovered as
above mentioned.

Since these problems cannot be solved in closed form, we will make use
of numerical optimization methods. The following section presents the
numerical results obtained for this problem.


% Considering the angles $\{\phi_i\}$ as the only design
% parameters, the following optimization problem can be formulated:
% \begin{equation}
%   \Phi^* = \min_{\Phi\in\mathcal{P}} J(\Phi)
% \end{equation}
% where $J$ is the cost function, $\Phi=[\phi_1\cdots\phi_6]^T$, and $\mathcal{P}$ is the set of
% allowable angle values. As cost function we consider a linear
% combination of the symmetric of the maximum force and torque,
% \begin{equation}
%   J(\Phi) = -(1-\lambda) F^\mathrm{max} - \lambda M^\mathrm{max}
% \end{equation}
% which corresponds to a linear scalarization of a multi-objective
% problem consisting of the objective functions~\eqref{eq:fmax}
% and~\eqref{eq:mmax}.  This problem is neither convex, smooth, not
% linear, so we will employ global optimization methods to solve it
% numerically. In particular, the results presented below were obtained
% using differential evolution algorithm~\cite{storn97}.

% Recall that the values of parameters $\{\theta_i\}$ and $\{w_i\}$ are
% set according to \eqref{eq:thetas} and \eqref{eq:ws}. For the
% $\{\phi_i\}$ we will consider three cases:
% \begin{enumerate}
% \item the simplest one consists in alternating a single angle value between
%   two symmetric values, similarly to what is done in~\cite{jiang13,voyles12}, that is
%   \begin{equation}
%     % \mathcal{P}_1 = \left\{ (\phi, -\phi, \phi, -\phi, \phi, -\phi)^T \::\:
%     %   \phi\in[0; \frac{\pi}{2}] \right\}
%     \mathcal{P}_1 = \left\{ \left(
%         \begin{array}{r}
%           \phi \\ -\phi \\ \phi \\ -\phi \\ \phi \\ -\phi
%         \end{array}
%       \right) \::\: |\phi|<\frac{\pi}{2}] \right\}
%   \end{equation}
% \item next we consider three degrees of freedom, corresponding to the
%   three principal axes joining pair os motors at opposite sides, that
%   is
%   \begin{equation}
%     % \mathcal{P}_3 = \left\{ (\phi_1, \phi_2, \phi_3, -\phi_1, -\phi_2, -\phi_3)^T \::\:
%     %   \phi_1\in[0; \frac{\pi}{2}], \phi_2,\phi_3\in[0; \frac{\pi}{2}] \right\}
%     \mathcal{P}_3 = \left\{ \left(
%         \begin{array}{r}
%           \phi_1 \\ \phi_2 \\ \phi_3 \\ -\phi_1 \\ -\phi_2 \\ -\phi_3
%         \end{array}
%       \right) \::\: |\phi_{1,\ldots,3}|<\frac{\pi}{2} \right\}
%   \end{equation}
% \item finally, there is the case where we have the six degrees of
%   freedom, one for each propeller, that
%   is
%   \begin{equation}
%     % \mathcal{P}_3 = \left\{ (\phi_1, \phi_2, \phi_3, -\phi_1, -\phi_2, -\phi_3)^T \::\:
%     %   \phi_1\in[0; \frac{\pi}{2}], \phi_2,\phi_3\in[0; \frac{\pi}{2}] \right\}
%     \mathcal{P}_3 = \left\{ \left(
%         \begin{array}{r}
%           \phi_1 \\ \phi_2 \\ \phi_3 \\ \phi_4 \\ \phi_5 \\ \phi_6
%         \end{array}
%       \right) \::\: |\phi_{1,\ldots,6}|<\frac{\pi}{2} \right\}
%   \end{equation}
% \end{enumerate}
% We have considered an unilateral interval for one of the propellers
% for symmetry breaking, since inverting the signs of all the $\phi_i$
% angles would result in the same cost function values.

\subsection{Numerical results}
\label{sec:numerical-results}

The optimization problem in~\eqref{eq:nbi} shows some features that
make it non-trivial to solve: it is both strongly non-convex with
mixed continuous and discrete variables. First, we will factor out the
discrete part by iteratively trying each combination of $\{w_i\}$
values modulo rotations (also called \emph{orbits}\footnote{For
  instance, $[1,-1,1,1,1,1]$ and $[1,1,-1,1,1,1]$ belong to the same
  orbit, and thus it is redundant to try both of them.}): from its
$2^6=64$ possible combinations, only 14 correspond to combinations
where no pair can be made equal after rotating one of them. Second, we
use a random multistart initialization together with a convex
optimization algorithm: for each sample drawn uniformly from the
$\{\phi_i\::\:|\phi_i| \leq \phi_{max}\}$ cube, we run the Constrained
Optimization BY Linear Approximation (COBYLA) algorithm~\cite{powell94},
as implemented in the SciPy optimization package.

% - explicar valores de K

To make the relation between force and actuation dimensionless, we
divided the actuation matrix by $K_1$. This way, the only dependence
on physical coefficients of this matrix is on the parameters $d$ and
the ratio $K_2/K_1$. Using~\eqref{eq:blade}, this ratio can be
expressed as
\begin{equation}
  \frac{K_2}{K_1} = \frac{D}{2\pi} \frac{C_P}{C_T}
\end{equation}
For various small propellers (of about 4'') we found\footnote{We used the
  UIUC Propeller Data Site, Vol. 2,
  http://m-selig.ae.illinois.edu/props/propDB.html (retrieved
  Nov-2015).} this ratio to be approximately~0.01. And thus we used
this value to obtain the numerical results presented in this section.
For $d$ we used the one from the design presented in
Section~\ref{sec:vehicle-design}, that is, $d=0.16$.

% - mostrar shadow optima

The extremes of the NBI search line are the shadow minima,
\textit{i.e.,} the minima of~\eqref{eq:minp}
and~\eqref{eq:minq}. For~1000 random initializations, we obtained
these values for the shadow minima: $(p^*,q^0)=(0.250, 15.01)$ and
$(p^0,q^*)=(0.4167,9.728)$.

% - mostrar resultados NBI

\begin{figure*}
  \centering
  \subfloat[]{\includegraphics[height=0.28\linewidth]{nbi_pq.png}} \hfill
  \subfloat[]{\includegraphics[height=0.28\linewidth]{nbi_fml.png}} \hfill
  \subfloat[]{\includegraphics[height=0.28\linewidth]{nbi_FM.png}}
  \caption{Pointwise approximation to the Pareto optimal set using the
    NBI method: (a)~obtained points in the $(p.q)$ space,
    (b)~$F^{max}$ and $M^{max}$ in function of $\lambda$, and (c)~in
    the $(F^{max},M^{max})$ space. The dimensions for $F^{max}$ and
    $M^{max}$ have no physical meaning because of the division of the
    actuation matrix by~$K_1$, as explained in the text, and thus are
    to be understood in relative terms only.}
  \label{fig:nbiout2}
\end{figure*}

\begin{table*}
  \centering
  \begin{tabular}{l|cccccc|rrrrrr|cc}
    $\lambda$
    & $\phi_1$ & $\phi_2$ & $\phi_3$ & $\phi_4$ & $\phi_5$ & $\phi_6$
    & $w_1$ & $w_2$ & $w_3$ & $w_4$ & $w_5$ & $w_6$
    & $F^{max}$ & $M^{max}$ \\
    \hline
    0
    & 54.74 &-54.73 & 54.74 &-54.74 & 54.73 &-54.74
    &-1 & 1 &-1 & 1 &-1 & 1
    & 2.000 & 0.2798 \\
    0.25
    & 54.73 &-54.74 & 54.74 &-54.73  & 54.74 &-54.74
    &-1 & 1 &-1 & 1 &-1 & 1
    & 2.000 & 0.2798 \\
    0.5
    & 53.73 &-53.73 & 53.73 &-53.73 & 53.73 &-53.73
    &-1 & 1 &-1 & 1 &-1 & 1
    & 1.999 & 0.2844 \\
    0.75
    & 49.56 &-49.56 & 49.56 &-49.56 & 49.56 &-49.56
    &-1 & 1 &-1 & 1 &-1 & 1
    & 1.969 & 0.3009 \\
    1
    & 38.84 &-38.85 & 38.84 &-38.84  & 38.85 &-38.84
    &-1 & 1 &-1 & 1 &-1 & 1
    & 1.745 & 0.3206 \\
    \hline
  \end{tabular}
  \caption{Some of the configurations obtained for 5 equally spaced
    values of $\lambda$. The values for $\{\phi_i\}$ are shown in
    degrees. As before, the dimensions for $F^{max}$ and
    $M^{max}$ have no physical meaning.}
  \label{tab:configs}
\end{table*}

With these values, we ran our optimization method for $\lambda$
ranging from~0 to~1 on 0.01 steps, for 1000 random initializations
each. The result is a set of points approximating the Pareto optimal
set, shown in Figure~\ref{fig:nbiout2}. The (a) plot of this figure
suggests a convex Pareto front. For a range of $\lambda$ values from 0
to about 0.4, the $(p,q)$ values are constant. As $\lambda$ increases
over 0.4, there is a drop on the values of $q$, meaning a slight
increase on $M^{max}$

% - justificar escolha

\begin{table}
  \centering
  \begin{tabular}{ccccccc}
    propeller ($i$) & 1 & 2 & 3 & 4 & 5 & 6 \\
    \hline
    $\theta_i$ & 0 & 60 & 120 & 180 & 240 & 300 \\
    $\phi_i$ & 55 & -55 & 55 & -55 & 55 & -55 \\
    $w_i$ & -1 & 1 & -1 & 1 & -1 & 1 \\
    \hline
  \end{tabular}
  \caption{Design parameters of the selected solution. Both
    $\{\theta_i\}$ and $\{\phi_i\}$ are expressed in degrees.}
  \label{tab:selected}
\end{table}

Table~\ref{tab:configs} shows some of the optimal configurations
obtained for some values of $\lambda$. In our choice we decided to
prefer maximum force,
since the corresponding maximum torque is not significantly
lower than other non-dominated solutions, and thus we selected a configuration found in
the lower range of $\lambda$ values. We rounded off the angle values
to the closest integer degree, resulting in the configuration shown in
Table~\ref{tab:selected}. All of the following results shown in this
paper employ this selected configuration.




% We applied this methodology to a vehicle with the design parameters
% shown in Table~\ref{tab:params}. The constants $K_1$ and $K_2$ have
% the same values than the ones reported in \cite{jiang13msc}, while $d$
% is the one from the design presented in
% Section~\ref{sec:specification}.

% \begin{table}
%   \centering
%   \begin{tabular}{ccccc}
%     $K_1$ & $K_2$ & $d$ & $\{\theta_i\}$  \\
%     \hline
%     5.7 & 1.3 & 0.16 & in \eqref{eq:thetas}  \\
%   \end{tabular}
%   \caption{Design parameters of the simulated vehicle.}
%   \label{tab:params}
% \end{table}

% To get an idea of the shape of the feasible set in the space of costs
% values, we sampled $N$ points uniformly from $\Psi$ and identified the
% non-dominated points. The result is shown in
% Figure~\ref{fig:pareto}. Although these non-dominated points are an
% unstable approximate to the Pareto optimal set, they suggest that a
% linear scalarization of the cost functions should provide a good
% approximation to the actual Pareto optimal set.

% \begin{figure}
%   \centering
%   \includegraphics[width=0.8\linewidth]{pareto_12.png}
%   \caption{Samples ($N=10^7$) uniformly drawn from the feasible set $\Psi$ shown
%     in the cost functions space. The non-dominated samples are shown
%     as bigger points.}
%   \label{fig:pareto}
% \end{figure}

% The shadow minimum is defined by the optimization of the individual
% cost functions, that is,
% \begin{equation}
%   J_1^* = \min_{\bar{\psi}} J_1(\bar{\psi}) \qquad
%   J_2^* = \min_{\bar{\psi}} J_2(\bar{\psi})
% \end{equation}
% After running differential evolution over~100 runs, we obtained the
% minimum values $J_1^*=-11.40$ and $J_2^*=-3.176$ (with a total variation
% below~$10^{-3}$).

% % The shadow minimum is defined by the optimization of the individual
% % cost functions, that is,
% % \begin{equation}
% %   \bar{\psi_1^*} = \argmin_{\bar{\psi}} J_1(\bar{\psi}) \qquad
% %   \bar{\psi_2^*} = \argmin_{\bar{\psi}} J_2(\bar{\psi})
% % \end{equation}
% % These minimizers corresponds to two points in the costs function
% % space, say $(J_1^*,J_2^+)=( J_1(\bar{\psi_1^*}),J_2(\bar{\psi_1^*}))$
% % and $(J_1^+,J_2^*)=( J_1(\bar{\psi_2^*}),J_2(\bar{\psi_2^*}))$. The
% % line between these two points is called the \emph{convex hull of
% %   individual minima} (CHIM)}. If the Pareto optimal set is convex,
% % which from Figure~\ref{fig:pareto} seems to be the case, then it
% % certainly lies below the CHIM line~\cite{das98}. After running
% % differential evolution over~100 runs, we obtained the mean values
% % $(J_1^*,J_2^+)=(-11.40, -2.606)$ and
% % $(J_1^+,J_2^*)=(-9.260 , -3.176)$.

% We applied the linear scalarization~\eqref{eq:linscal} and used a
% differential evolution\footnote{We used the implementation found in
%   the \emph{optimize} package of SciPy. We used the default parameters
%   except for a population size of 512 and a mutation dithering range
%   of~$[0.2; 0.6]$.} algorithm~\cite{storn97} on this cost function,
% while ranging $\lambda$ over an equally spaced grid from 0 to 1
% totalling 20 runs. Figure~\ref{fig:pm12} shows the obtained
% results. The bottom left solution in the \ref{fig:pm12}(b) plot, which
% approximately dominates all of the others, includes 12 of these 20
% runs. 

%  \begin{figure}
%   \centering
%   % data from rm12dep_2.npy
%   \subfloat[]{\includegraphics[width=0.33\linewidth]{out_pm12_cost.pdf}}
%   \subfloat[]{\includegraphics[width=0.33\linewidth]{out_pm12_pareto.pdf}}
%   \subfloat[]{\includegraphics[width=0.33\linewidth]{out_pm12_solution.pdf}}
%   \caption{Optimization results,
%     from left to right: (a)~plot of $J$, $J_1=-F^\mathrm{max}$, and
%     $J_2=M^\mathrm{max}$ in function of $\lambda$, (b)~solutions found
%     shown in the $J_1$ and $J_2$ space, and (c)~optimized $\{\phi_i\}$
%     in function of~$\lambda$.}
%   \label{fig:pm12}
% \end{figure}

% We then ran differential evolution 100 times with
% $\lambda=0.4$, an heuristically chosen within a range of $\lambda$
% values where the resulting costs seem stable. From these runs, the
% dominant solution, found in 24 of these 100 runs, is shown in
% Table~\ref{tab:best}. Its costs are $J_1=-11.40$ and $J_2=-3.176$,
% which coincide with the individual minima $J_1^*$ and $J_2^*$ found
% above. That is, no other solution found has lower cost in either of
% the cost functions. 

% \begin{table}
%   \centering
%   \begin{tabular}{ccccccc}
%     propeller ($i$) & 1 & 2 & 3 & 4 & 5 & 6 \\
%     \hline
%     $\theta_i$ & 0\textdegree & 60\textdegree & 120\textdegree & 180\textdegree & 240\textdegree & 300\textdegree \\
%     $\phi_i$ & -74.54\textdegree & -39.54\textdegree & -74.57\textdegree & -39.48\textdegree & -74.49 \textdegree & -39.46 \\
%     $w_i$ & -1 & 1 & -1 & 1 & -1 & 1 \\
%     \hline
%   \end{tabular}
%   \caption{The solution found over 100 runs of the optimizer with
%     $\lambda=0.4$.}
%   \label{tab:best}
% \end{table}

% To wrap up these results, Figure~\ref{fig:costs} shows the solutions
% found by differential evolution as boxplots, in three different
% scenarios: the solutions for the two shadow minima
% $\bar{\psi}_1^*=\arg\min J_1(\bar{\psi})$ and
% $\bar{\psi}_2^*=\arg\min J_2(\bar{\psi})$, and the solutions
% $\bar{\psi}_\lambda^*=\arg\min J(\bar{\psi}; \lambda)$ for
% $\lambda=0.4$. Note that the minimum costs of the
% $\bar{\psi}_\lambda^*$ solutions coincide with the individual
% minima of $J_1^*=J_1(\bar{\psi}_1^*)$ and $J_2^*=J_2(\bar{\psi}_2^*)$.

% \begin{figure}
%   \centering
%   \includegraphics[width=0.8\linewidth]{out_osom12.pdf}
%   \caption{Boxplots of the $J_1$ and $J_2$ costs found while
%     separately minimizing, $J_1(\bar{\psi})$ (first and second
%     boxplot), $J_2(\bar{\psi})$ (third and fourth boxplot), and
%     $J(\bar{\psi}; \lambda)$ for $\lambda=0.4$ (fifth and
%     sixth). Each boxplots represents the minimum, maximum, and median
%     values.}
%   \label{fig:costs}
% \end{figure}

% The solution in Table~\ref{tab:best} implies three of the propellers to be
% almost horizontal, hindering the air flow to produce the desired
% thrust. Therefore, we decided on a different solution, with the angle
% values rounded from an handpicked solution from the Pareto optimal set
% obtained above. Its parameters are shown in Table~\ref{tab:selected}
% and its costs are $J_1=-F^\mathrm{max}=-11.40$ and
% $J_2=-M^\mathrm{max}=-2.623$, that is, the same maximum force but with
% 82\% of the maximum torque. We preferred this solution
% because the propeller angles are not as horizontal, resulting in a
% significantly better air flow.
% % Running the
% % optimization algorithm on a restricted feasibility space parametrized
% % by~$\phi$ defined as
% % \begin{equation}
% %   \bar{\psi} = \left(-\phi, \phi, -\phi, \phi, -\phi, \phi, -1, -1, 1, 1,
% %     1, -1\right)^T
% % \end{equation}
% % with $\lambda=0.4$ resulted 

% \begin{table}
%   \centering
%   \begin{tabular}{ccccccc}
%     propeller ($i$) & 1 & 2 & 3 & 4 & 5 & 6 \\
%     \hline
%     $\theta_i$ & 0\textdegree & 60\textdegree & 120\textdegree & 180\textdegree & 240\textdegree & 300\textdegree \\
%     $\phi_i$ & -55\textdegree & 55\textdegree & -55\textdegree & 55\textdegree & -55\textdegree & 55\textdegree \\
%     $w_i$ & -1 & -1 & 1 & 1 & 1 & -1 \\
%     \hline
%   \end{tabular}
%   \caption{Design parameters of the selected solution.}
%   \label{tab:selected}
% \end{table}


% Figure~\ref{fig:forces} shows the maximum forces and torques in set
% $\mathcal{P}_1$. The cost functions~\eqref{eq:fmax}
% and~\eqref{eq:mmax} correspond to the symmetric of the pointwise
% minimum over the three lines for each one of these plots. The
% non-convexity of the cost functions is clearly visible in this plot. When
% $\phi=0$ the maximum force is zero, since in this configuration all
% thrust is directed along the $Z$ axis: this corresponds to the
% usual configuration of hexacopters (and quadrotors, although with less
% propellers). As $|\phi|$ increases, the maximum force along directions
% other than $Z$ increases.

% \begin{figure}
%   \centering
%   \includegraphics[width=\linewidth]{out_plot_forces.pdf}
%   \caption{Maximum force and torque, $F^\mathrm{max}$ and
%     $M^\mathrm{max}$, for each one of
%     the body frame axes while ranging across the design space
%     $\mathcal{P}_1$.}
%   \label{fig:forces}
% \end{figure}


% \begin{figure}
%   \centering
%   \includegraphics[width=0.32\linewidth]{out_pm1_cost.pdf} \hfill
%   \includegraphics[width=0.32\linewidth]{out_pm1_curve.pdf} \hfill
%   \includegraphics[width=0.32\linewidth]{out_pm1_solution.pdf} 
%   \caption{Optimization results for feasibility set $\mathcal{P}_1$,
%     from left to right: (a)~plot of $-J$, $F^\mathrm{max}$, and
%     $M^\mathrm{max}$ in function of $\lambda$, (b)~relationship between $F^\mathrm{max}$ and
%     $M^\mathrm{max}$, and (c)~optimized $\phi$ in function of $\lambda$.}
%   \label{fig:pm1}
% \end{figure}

% \begin{figure}
%   \centering
%   \includegraphics[width=0.32\linewidth]{out_pm3_cost.pdf} \hfill
%   \includegraphics[width=0.32\linewidth]{out_pm3_curve.pdf} \hfill
%   \includegraphics[width=0.32\linewidth]{out_pm3_solution.pdf} 
%   \caption{Optimization results for feasibility set $\mathcal{P}_3$,
%     same description as in Figure~\ref{fig:pm1}}
%   \label{fig:pm3}
% \end{figure}

%%% Local Variables:
%%% mode: latex
%%% TeX-master: "main"
%%% End:
